% Shared preamble of the HSEMHS mode documents (docs/physics/*.tex, docs/technical/*.tex).
\documentclass[10pt]{extarticle}
\usepackage[a4paper,margin=1.7cm]{geometry}
\usepackage[T1]{fontenc}\usepackage{lmodern}\usepackage{tabularx}
\usepackage{amsmath,amssymb,bm,enumitem,titlesec,xcolor,booktabs,array,tikz,graphicx,longtable}
\usepackage[hidelinks]{hyperref}
\titleformat{\section}{\large\bfseries\color{blue!50!black}}{\thesection.}{0.5em}{}
\titlespacing*{\section}{0pt}{9pt}{3pt}
\titleformat{\subsection}{\normalsize\bfseries}{\thesubsection}{0.5em}{}
\titlespacing*{\subsection}{0pt}{6pt}{2pt}
\setlist{nosep,leftmargin=1.3em}
\setlength{\parindent}{0pt}\setlength{\parskip}{4pt}
\renewcommand{\arraystretch}{1.15}
\newcommand{\B}{\mathbf B}\newcommand{\J}{\mathbf J}\newcommand{\F}{\mathbf F}\newcommand{\R}{\mathbf R}
\newcommand{\zh}{\hat{\mathbf z}}\newcommand{\Bo}{\mathbf B_{\rm obs}}\newcommand{\Lv}{\mathbf L}
\newcommand{\code}[1]{\texttt{#1}}
\newcolumntype{L}[1]{>{\raggedright\arraybackslash}p{#1}}
\newcolumntype{Y}{>{\raggedright\arraybackslash}X}
\newcommand{\doctitle}[2]{{\Large\bfseries #1}\\[2pt]{\small #2}\par\vspace{2pt}}
\pagestyle{plain}
\AtBeginDocument{\sloppy}

\begin{document}
\doctitle{mhs-poisson-ltau: pressure-only magnetohydrostatic correction ($\tau$ grid)}{Physics of \code{run.mode = "mhs-poisson-ltau"} (HSEMHS). 24 Oct 2026.}

\section{What the mode does}
An inversion delivers $T$ and $\B$ on surfaces of constant $\log_{10}\tau_{500}$. To use them physically we need the gas
pressure, the density and the \textbf{geometric height} of every $\tau$ level, i.e.\ where in space the surfaces lie.
\code{column-hse-ltau} answers this column by column. This mode adds the horizontal force balance between columns and
the Lorentz force of the given field: with $T(\tau)$ and $\B(\tau)$ fixed, it finds $\ln P(\tau)$ and the heights $z(\tau)$
that minimise the force imbalance, in one non-linear least-squares solve directly on the $\tau$ grid. It is the
optical-depth counterpart of \code{mhs-poisson-z}, and the mode intended for inversion output.

\section{Quantities}
\begin{tabular}{>{$}l<{$} p{11.2cm} l}
\toprule
\text{Symbol} & Meaning & Unit \\ \midrule
\tau & optical depth at 500\,nm; levels uniform in $\log_{10}\tau$ & -- \\
T,\ \B & temperature and field on the $\tau$ levels, fixed & K, T \\
P,\ u & gas pressure on the $\tau$ levels; $u=\ln P$ & Pa, -- \\
z(\tau) & height of a level in a column; follows from $P$ (Sect.~3) & m \\
c(x,y) & height of the cut level $\tau_c$ in a column (one unknown per column) & m \\
\chi,\ \rho & extinction and density from the EOS, functions of $(T,P)$ & m$^{-1}$, kg\,m$^{-3}$ \\
\Lv & Lorentz force $\J\times\B$, $\J=\nabla\times\B/\mu_0$ in physical space & N\,m$^{-3}$ \\
\tau_c & cut: levels with $\log_{10}\tau\ge\log_{10}\tau_c$ are solved (production $-4$, $z\approx0.56$\,Mm) & -- \\
\bottomrule
\end{tabular}

\section{Heights follow from the pressure}
On a $\tau$ grid the geometry is not given: the spacing of two levels is $dz=-d\tau/\chi$, and $\chi$ depends on $P$.
So the heights are part of the solution. In each column
\[
z_k=c+\sum_{m=k_c}^{k-1}\Delta z_m,\qquad \Delta z_m=-\int_{\ln\tau_m}^{\ln\tau_{m+1}}\frac{\tau}{\chi(T,P)}\,d\ln\tau ,
\]
i.e.\ the height of every solved level is the height $c$ of the cut level plus the thicknesses of the levels above it,
each set by its own opacity. A change of pressure at one level shifts every level below it. Above the cut the column is
kept hydrostatic in $\tau$ (as in \code{column-hse-ltau}) and moves rigidly with $c$.

\section{The force balance in $\tau$ coordinates}
The static balance $\F=-\nabla P+\Lv-\rho g\zh=0$ must hold in physical space, but our derivatives are along the $\tau$
surfaces, which are tilted. For any quantity, $\partial_x f|_z=\partial_x f|_\tau-\partial_xz|_\tau\,\partial_zf$. Using the
vertical balance $\partial_zP=-(\rho g-L_z)$:
\[
\boxed{\;F_x=-\partial_xP\big|_\tau-(\rho g-L_z)\,\partial_xz\big|_\tau+L_x,\qquad
\frac{d\ln P}{d\ln\tau}=\frac{(\rho g-L_z)\,\tau}{\chi\,P}\;}
\]
(and the same for $y$). The horizontal equation is the familiar one with a geometric term: on a tilted $\tau$ surface,
part of the weight of the gas acts sideways. The vertical equation is hydrostatic equilibrium in $\tau$ with the
vertical Lorentz force reducing the effective weight. Both are divided by $P$ (a relative imbalance, as in
\code{mhs-poisson-z}), and the objective
\[
\Phi=\sum_{\text{levels}}w\,\bar H^2\big(R_x^2+R_y^2\big)+\sum_{\text{intervals}}\frac{\bar H^2}{\Delta\bar z}R_z^2
\]
weights each level by its thickness and median scale height, so $R\bar H\approx F/(\rho g)$ layer by layer.

$\Lv$ is computed from $\B$ on the $\tau$ surfaces at their current heights. When the heights change, the vertical
spacing of the field samples changes and so do the currents; the solver follows this (Sect.~6 of the technical
document). The field values themselves are never changed.

\section{What decides the heights}
Each column's heights relative to its neighbours are set by the horizontal balance: a column whose gas is lighter at a
given height must be compressed (higher pressure) or displaced until the sideways pressure push vanishes. In a
strong-field column, the magnetic pressure carries part of the load, so less gas pressure is needed, the gas is less
dense and less opaque, and $\tau=1$ lies deeper: the \textbf{Wilson depression} emerges from the force balance, not
from an assumption. One global constant (all columns moved together) is not fixed by the equations; it is set by
requiring the mean height of $\tau=1$ to be zero.

\section{Results on the production cube}
\code{IO/input-mhs-ltau.nc}, $504\times504\times95$ levels; 17 Gauss--Newton steps, 14\,min on one node (8.3\,min on four).
Force imbalance after the solve, RMS over log-$\tau$ bands in units of $\rho g$: horizontal 0.04--0.09 and vertical
0.04--0.06 for $\log\tau<0$, 0.17 / 0.06 below $\tau=1$ (start: 0.3--0.8). $z(\tau=1)$ relative to the cube mean, median
[5\,\%, 95\,\%] in km, binned by $|\B|$ at $\tau=1$:
\begin{center}\small
\begin{tabular}{lccc}
\toprule
$|\B|$ at $\tau=1$ & column HSE in $\tau$ & fixed-point scheme (superseded) & \textbf{mhs-poisson-ltau} \\ \midrule
0--200\,G (245\,123 columns) & $-38$ [$-85$, 187] & $-55$ [$-257$, 380] & $-10$ [$-35$, 73] \\
200--500\,G & $-68$ [$-88$, 58] & $+120$ [$-180$, 549] & $-26$ [$-52$, 24] \\
500--1000\,G & $-86$ [$-104$, $-36$] & $+183$ [$-198$, 587] & $-51$ [$-89$, $-17$] \\
1000--1500\,G & $-121$ [$-144$, $-92$] & $+314$ [$-117$, 563] & $-96$ [$-125$, $-63$] \\
1500--3000\,G (440) & $-156$ [$-175$, $-140$] & $+402$ [58, 568] & $-130$ [$-167$, $-104$] \\
\midrule
columns outside $-200\ldots+300$\,km & 2290 & 69\,522 & 77 \\
\bottomrule
\end{tabular}
\end{center}
\begin{itemize}
\item The depression grows monotonically with field strength, somewhat less than in independent columns, because the
neighbours now share the load through the horizontal balance.
\item In weak-field regions the column-to-column scatter of $\tau=1$ shrinks by a factor of about 3: horizontal pressure
balance couples the columns that column HSE treats separately.
\item An earlier scheme alternated a height-grid Poisson solve with re-integration of the $\tau$ columns. It lifted the
strong-field columns instead of depressing them: it integrated the heights with a pressure that it then replaced. It is
kept only as history (\code{scripts/poisson/poisson\_tau.jl}).
\end{itemize}

\section{Limitations}
\begin{itemize}
\item As in \code{mhs-poisson-z}, a fixed $\B$ cannot balance the buoyancy of the temperature structure; the remaining
imbalance is of that kind. Above the cut ($\log\tau<-4$) the columns are only hydrostatic.
\item The heights of a few strong-field columns are only weakly constrained: when the solve stops, they still move by a
few km per step while $\Phi$ changes by less than $10^{-4}$. Different runs agree to about 1\,m at the median and
12--24\,m at the 99th percentile.
\item The vertical integration uses the trapezoid rule on the 0.1\,dex grid; sharp temperature jumps on that grid make
individual level thicknesses uncertain (see the cap in \code{column-hse-ltau}).
\end{itemize}
\end{document}
